\documentclass[11pt]{article}
\usepackage[T1]{fontenc}
\usepackage{lmodern}
\usepackage[margin=1in]{geometry}
\usepackage{amsmath,amssymb,amsthm,mathtools,booktabs,array}
\usepackage[hidelinks]{hyperref}
\hypersetup{pdftitle={Translation relations and a twist-free converse theorem at level 32},
 pdfsubject={Generalized translation lemmas and the level 32 converse theorem}}
\newtheorem{theorem}{Theorem}
\newtheorem{proposition}[theorem]{Proposition}
\newtheorem{lemma}[theorem]{Lemma}
\newtheorem{corollary}[theorem]{Corollary}
\theoremstyle{remark}
\newtheorem{remark}[theorem]{Remark}
\newcommand{\HH}{\mathfrak H}
\newcommand{\SL}{\mathrm{SL}}
\newcommand{\eps}{\varepsilon}
\newcommand{\bb}{\beta}
\newcommand{\tr}{\operatorname{tr}}
\title{Translation relations and a twist-free converse theorem at level $32$}
\author{}
\date{28 October 2026}

\begin{document}
\maketitle
\begin{abstract}
We prove the holomorphic Conrey--Farmer converse theorem at level $32$
for every positive even weight and both Fricke eigenvalues. Only the
Euler conditions at $2$ and $3$ are needed. The proof develops a signed
two-term elliptic criterion, applies it uniformly to the rescaled
$m=3$ relations at levels $4,8,16,32$, and completes the level $32$
case by an explicit character of $\Gamma_0(8)$. The character is
nontrivial on every cusp stabilizer, which proves cuspidality without
additional coefficient bounds. We also record a parameterized version
of the three-dimensional translation module used at level $36$.
This is a separate continuation of the level $36$ work; that work is
not modified here.
\end{abstract}

\section{The new level and its hypotheses}

The transformation-law results of Conrey and Farmer \cite{CF};
Conrey, Farmer, Odgers, and Snaith \cite{CFOS}; and
Bedert, Cooper, Oliver, and Zhang \cite{BCOZ} cover, collectively,
$1\leq N\leq18$ and $N=20,23,24$. Level $32$ is outside these lists.
The assertion below includes a separate proof of vanishing at all cusps;
we do not infer cuspidality from the functional equation alone.

\begin{theorem}\label{thm:32}
Let $k\geq2$ be even and let $a_1=1$. Suppose
\[
 F(s)=\sum_{n\geq1}a_n n^{-s}
\]
converges in some right half-plane and satisfies the following conditions:
\begin{enumerate}
\item The function
\[
 \Lambda(s)=\left(\frac{\sqrt{32}}{2\pi}\right)^s\Gamma(s)F(s)
\]
extends to an entire function bounded in every closed vertical strip,
and $\Lambda(s)=\omega\Lambda(k-s)$ for some $\omega\in\{1,-1\}$.
\item $a_{2n}=0$ for every $n\geq1$.
\item For every $n\geq1$,
\[
 a_{3n}=a_3a_n-3^{k-1}a_{n/3},
 \qquad a_{n/3}=0\quad\hbox{if }3\nmid n.
\]
\end{enumerate}
Then
\[
 f(z)=\sum_{n\geq1}a_n e^{2\pi inz}\ \in S_k(\Gamma_0(32)).
\]
Both functional-equation signs are allowed, and there is no additional
growth assumption on the coefficients.
\end{theorem}

Conditions (2)--(3) are precisely the coefficient conditions imposed by
the partial Euler product
\[
 F(s)=\frac{1}{1-a_3 3^{-s}+3^{k-1-2s}}
       \sum_{(n,6)=1}a_n n^{-s}.
\]
Consequently the full Conrey--Farmer Euler product at level $32$
implies the local hypotheses of the theorem. No local factor at any
prime other than $2$ and $3$ is used.

For $M=\left(\begin{smallmatrix}a&b\\c&d\end{smallmatrix}\right)
\in\mathrm{GL}_2^+(\mathbb R)$ use the normalized slash action
\[
 (v|_kM)(z)=(\det M)^{k/2}(cz+d)^{-k}
 v\left(\frac{az+b}{cz+d}\right).
\]
We suppress the subscript. Since $k$ is even, all nonzero scalar
matrices act trivially. Put
\[
 T=\bb(1),\qquad
 \bb(t)=\begin{pmatrix}1&t\\0&1\end{pmatrix},\qquad
 S=\begin{pmatrix}0&-1\\1&0\end{pmatrix},
\]
\[
 H_N=\begin{pmatrix}0&-1\\N&0\end{pmatrix},\qquad
 W_M=\begin{pmatrix}1&0\\M&1\end{pmatrix}.
\]
Convergence at a real point implies $a_n=O(n^A)$ for some $A\geq0$.
Thus $f$ is holomorphic on $\HH$. Mellin inversion and the functional
equation give
\begin{equation}\label{eq:fricke}
 f|H_N=\eps f,\qquad \eps=(-1)^{k/2}\omega\in\{1,-1\}.
\end{equation}
This is the usual Hecke argument, with the analytic root number
$\omega$ distinguished from the slash eigenvalue $\eps$.

\section{Two generalizations of the level \texorpdfstring{$36$}{36} method}

\subsection{The translation module with a parameter}

\begin{proposition}\label{prop:module}
Suppose $h$ is holomorphic on $\HH$, $h|S=\eps h$ with $\eps^2=1$,
and, for some $\tau\in\mathbb C$,
\[
 h|(T^2-\tau T+1)=0.
\]
Set $g=h|T$ and $v=g|S$. The span of $h,g,v$ is stable under
$\SL_2(\mathbb Z)$, with action
\begin{equation}\label{eq:module}
\begin{array}{c|cc}
 & |S & |T\\ \hline
 h & \eps h & g\\
 g & v & \tau g-h\\
 v & g & \tau h-\eps v
\end{array}
\end{equation}
No linear independence is required. The formal matrices are
\[
 A_\eps=\begin{pmatrix}\eps&0&0\\0&0&1\\0&1&0\end{pmatrix},
 \qquad
 B_{\tau,\eps}=\begin{pmatrix}0&1&0\\-1&\tau&0\\\tau&0&-\eps\end{pmatrix},
\]
and they satisfy
\[
 A_\eps^2=(A_\eps B_{\tau,\eps})^3=I,
 \qquad
 \det(XI-B_{\tau,\eps})=(X+\eps)(X^2-\tau X+1).
\]
\end{proposition}
\begin{proof}
The translation equation gives $h|T^{-1}=\tau h-g$.
The identity $STS=T^{-1}ST^{-1}$, projectively, gives
\[
 \eps v=h|STS=(\tau\eps h-v)|T^{-1}.
\]
Right multiplication by $T$ yields $v|T=\tau h-\eps v$.
All the other entries in the table follow directly from the definitions.
The matrix identities follow by multiplication.
\end{proof}

The values $\tau=-1,0,1$ correspond to Fourier frequencies with
translation eigenvalues that are primitive third, fourth, and sixth
roots, respectively. Thus the same construction applies after rescaling
at levels $9,16,36$. The module itself does not establish cuspidality
or the required congruence kernel in every case; these remain separate
steps. In particular, the Eisenstein caveat at levels $9$ and $16$ in
\cite[Section~5]{CF} must be respected.

\subsection{A signed two-term elliptic criterion}

\begin{lemma}[Elliptic vanishing]\label{lem:vanish}
Let $k>0$ be even and $u$ holomorphic on $\HH$. If
$u|E=\xi u$, where $E\in\SL_2(\mathbb R)$ is elliptic of infinite
order and $\xi$ is a root of unity, then $u=0$.
\end{lemma}
\begin{proof}
Taking a positive power reduces to $u|E=u$, with $E$ still elliptic
of infinite order. If $a\in\HH$ is its fixed point, the coordinate
$w=(z-a)/(z-\bar a)$ identifies its action with $w\mapsto\zeta w$,
where $\zeta$ has infinite order on the unit circle. Write
$u(z)(dz)^{k/2}=U(w)(dw)^{k/2}$. Invariance gives
\[
 U(w)=\zeta^{k/2}U(\zeta w).
\]
For $U(w)=\sum_{j\geq0}c_jw^j$ it follows that
$(1-\zeta^{j+k/2})c_j=0$. Since $j+k/2>0$, all $c_j$ vanish.
\end{proof}

\begin{lemma}[Symmetric pair]\label{lem:pair}
Suppose $h|T=\eta h$ with $\eta$ a root of unity. Let
$C\in\SL_2(\mathbb R)$, $m\in\mathbb Z$, $t\in\mathbb R$, and
$\lambda\in\{1,-1\}$. Put $C'=T^mC^{-1}T^{-m}$.
If
\begin{equation}\label{eq:pair}
 h|(1-\lambda C)\bb(-t)+h|(1-\lambda C')\bb(t)=0
\end{equation}
and
\[
 E=C^{-1}\bb(2t-m)
\]
is elliptic of infinite order, then $h|C=\lambda h$.
\end{lemma}
\begin{proof}
Let $u=h|(1-\lambda C)$. Using only the known scalar action of
$T^m$, we obtain
\[
 h|(1-\lambda C')
 =\eta^m h|(1-\lambda C^{-1})T^{-m}
 =-\lambda\eta^m u|C^{-1}T^{-m}.
\]
Thus \eqref{eq:pair}, right-multiplied by $\bb(t)$, becomes
\[
 u=\lambda\eta^m u|C^{-1}\bb(2t-m).
\]
Apply Lemma~\ref{lem:vanish}.
\end{proof}

This formulation permits signed transformation laws and integer
translation conjugacies. It is useful to check the trace directly:
if $C=\left(\begin{smallmatrix}a&b\\c&d\end{smallmatrix}\right)$, then
\[
 \tr\bigl(C^{-1}\bb(s)\bigr)=a+d-cs.
\]
A determinant-one real matrix with rational, nonintegral trace in
$(-2,2)$ is elliptic of infinite order. Indeed, finite order would make
its trace a rational algebraic integer, hence an integer.

\section{The uniform rescaled relation at levels \texorpdfstring{$4,8,16,32$}{4,8,16,32}}

We first record the needed consequence of a good Euler factor, including
the justification of the left replacements in the group algebra.

\begin{lemma}[Prime relation]\label{lem:prime}
Suppose $f|T=f$, $f|H_N=\eps f$, $\eps^2=1$, and for a prime
$p\nmid N$ the coefficients of $f$ satisfy
$a_{pn}=a_pa_n-p^{k-1}a_{n/p}$. For each nonzero residue class $b$
modulo $p$, choose an integer $c$ such that $Nbc\equiv-1\pmod p$,
and put
\[
 \gamma_b=\begin{pmatrix}p&-b\\-Nc&(1+Nbc)/p\end{pmatrix}.
\]
Then
\begin{equation}\label{eq:prime}
 f\bigg|\sum_{b\bmod p}^{*}(1-\gamma_b)\bb(b/p)=0.
\end{equation}
\end{lemma}
\begin{proof}
The coefficient recurrence is equivalent to
\[
 f|\mathcal T_p=p^{1-k/2}a_pf,\qquad
 \mathcal T_p=\begin{pmatrix}p&0\\0&1\end{pmatrix}
 +\sum_{r=0}^{p-1}\begin{pmatrix}1&r\\0&p\end{pmatrix}.
\]
Conjugate by $H_N$ and subtract. The two diagonal terms cancel.
Right multiplication by $\operatorname{diag}(p,1)$ gives
\[
 f\bigg|\sum_{r=1}^{p-1}\bb(r/p)
 =f\bigg|\sum_{c=1}^{p-1}\begin{pmatrix}p^2&0\\-Ncp&1\end{pmatrix}.
\]
Periodicity permits any representatives for $r$. Also
$f|W_N=f$, since $W_N=H_NT^{-1}H_N^{-1}$.
Replacing $c$ by $c+pj$ on the right is therefore permitted:
it amounts to left multiplication by $W_N^{-j}$.
Reindex by $Nbc\equiv-1\pmod p$ and use the exact identity
\[
 \gamma_b\bb(b/p)=\begin{pmatrix}p&0\\-Nc&1/p\end{pmatrix}.
\]
Scalar matrices do not affect the slash action, giving \eqref{eq:prime}.
This is the prime case of \cite[Corollary~2]{CF}.
\end{proof}

Now let $N=4M$ with $3\nmid M$. Suppose $f$ has the functional
equation at level $N$, $a_{2n}=0$, and the good Euler recurrence at $3$.
Put $D=\operatorname{diag}(1,2)$ and
\[
 h(z)=f(z/2)=2^{k/2}(f|D)(z).
\]
Odd Fourier support and Fricke conjugation give
\begin{equation}\label{eq:half}
 h|T=-h,\qquad h|H_M=\eps h,\qquad h|W_M=-h.
\end{equation}
The last equality uses $W_M=H_MT^{-1}H_M^{-1}$ and $\eps^2=1$.

Define $\sigma\in\{1,-1\}$ by $\sigma M\equiv-1\pmod3$.
In Lemma~\ref{lem:prime}, choose $b=-1,c=-\sigma$ and
$b=1,c=\sigma$. After conjugating by $D$, the relation is
\begin{equation}\label{eq:AB}
 h|(1-A)\bb(-2/3)+h|(1-B)\bb(2/3)=0,
\end{equation}
where
\[
 A=\begin{pmatrix}3&2\\2\sigma M&(1+4\sigma M)/3\end{pmatrix},
 \qquad
 B=\begin{pmatrix}3&-2\\-2\sigma M&(1+4\sigma M)/3\end{pmatrix}.
\]
Set
\begin{equation}\label{eq:Cgeneral}
 C=-W_M^{-\sigma}A
  =\begin{pmatrix}-3&-2\\\sigma M&(2\sigma M-1)/3\end{pmatrix},
 \qquad
 C'=-W_M^{\sigma}B
  =\begin{pmatrix}-3&2\\-\sigma M&(2\sigma M-1)/3\end{pmatrix}.
\end{equation}
All four matrices have determinant one. Because $h|W_M^{\pm1}=-h$,
\eqref{eq:AB} becomes
\begin{equation}\label{eq:CC}
 h|(1+C)\bb(-2/3)+h|(1+C')\bb(2/3)=0.
\end{equation}
A direct multiplication gives
\begin{equation}\label{eq:r}
 C'=\bb(r)C^{-1}\bb(-r),\qquad
 r=\frac23+\frac{8}{3\sigma M}.
\end{equation}

\begin{proposition}\label{prop:four}
Under the assumptions just made, if $M\in\{1,2,4,8\}$, then
$h|C=-h$ for the matrix $C$ in \eqref{eq:Cgeneral}.
\end{proposition}
\begin{proof}
For these four values, $r$ is an integer. Apply Lemma~\ref{lem:pair}
with $\eta=-1$, $\lambda=-1$, $m=r$, and $t=2/3$.
The determinant-one matrix
\[
 E=C^{-1}\bb(4/3-r)
\]
has trace
\[
 \frac{2\sigma M-10}{3}
 -\sigma M\left(\frac23-\frac{8}{3\sigma M}\right)
 =-\frac23.
\]
It is therefore elliptic of infinite order, proving the assertion.
\end{proof}

\begin{center}
\begin{tabular}{@{}rrrrr@{}}
\toprule
$M$ & $N=4M$ & $\sigma$ & $r$ & $\tr E$\\
\midrule
1 & 4  & $-1$ & $-2$ & $-2/3$\\
2 & 8  & $ 1$ & $ 2$ & $-2/3$\\
4 & 16 & $-1$ & $ 0$ & $-2/3$\\
8 & 32 & $ 1$ & $ 1$ & $-2/3$\\
\bottomrule
\end{tabular}
\end{center}

The range is exact for this particular integer-translation argument:
if $r$ in \eqref{eq:r} is integral, then
$\sigma M(3r-2)=8$, so $M$ divides $8$. For other $M$, the same
matrix identities remain valid, but \eqref{eq:half} does not provide
a scalar action of the nonintegral translation $\bb(r)$.
No such scalar action is assumed here.

\section{Completion of the level \texorpdfstring{$32$}{32} transformation law}

Specialize to $M=8$. We have
\[
 A=\begin{pmatrix}3&2\\16&11\end{pmatrix},\qquad
 B=\begin{pmatrix}3&-2\\-16&11\end{pmatrix},\qquad
 C=\begin{pmatrix}-3&-2\\8&5\end{pmatrix},\qquad
 C'=\begin{pmatrix}-3&2\\-8&5\end{pmatrix}.
\]
For clarity, the complete defect calculation in this case is
\[
 W_8^{-1}A=-C,\qquad W_8B=-C',\qquad C'=TC^{-1}T^{-1}.
\]
Let $u=h|(1+C)$. Since $h|T=-h$,
\[
 h|(1+C')=-u|C^{-1}T^{-1}.
\]
Equation \eqref{eq:CC} gives
\begin{equation}\label{eq:explicitE}
 u=u|E,\qquad
 E=C^{-1}\bb(1/3)
  =\begin{pmatrix}5&11/3\\-8&-17/3\end{pmatrix},
 \qquad \tr E=-2/3.
\end{equation}
Thus $u=0$ and
\begin{equation}\label{eq:gensigns}
 h|T=h|W_8=h|C=-h.
\end{equation}

\begin{lemma}\label{lem:generators}
Projectively,
\[
 \Gamma_0(8)=\langle T,W_8,C\rangle.
\]
\end{lemma}
\begin{proof}
First, projectively, $\Gamma(2)$ is generated by
$T^2$ and $\left(\begin{smallmatrix}1&0\\2&1\end{smallmatrix}\right)$.
Here is the elementary reduction argument. In the first column of an
element of $\Gamma(2)$, the top entry is odd and the bottom entry is
even. Left multiplication by powers of these two generators subtracts
even multiples of one entry from the other. Choosing a nearest multiple
makes the new absolute value strictly smaller than the other one:
equality is impossible because of the different parities. Repeating
reduces the bottom entry to zero, after which the matrix is a power of
$T^2$, up to sign.

Conjugation by $D=\operatorname{diag}(1,2)$ identifies $\Gamma(2)$
with $\Gamma_0(4)$ and carries these generators to $T,W_4$.
The homomorphism
\[
 \Gamma_0(4)\longrightarrow\mathbb Z/2\mathbb Z,
 \qquad
 \begin{pmatrix}a&b\\4c&d\end{pmatrix}\longmapsto c\pmod2
\]
has kernel $\Gamma_0(8)$. Using representatives $1,W_4$, Schreier's
lemma gives generators
\[
 T,\qquad W_4^2=W_8,\qquad W_4TW_4^{-1}.
\]
Finally,
\[
 C=(W_4^{-1}T)^2,\qquad
 W_4TW_4^{-1}=W_8CT^{-1}.
\]
This proves the assertion, and also verifies directly the generating
set displayed in \cite[Section~2]{CF}.
\end{proof}

For $\gamma=\left(\begin{smallmatrix}a&b\\8c&d\end{smallmatrix}\right)
\in\Gamma_0(8)$ define
\begin{equation}\label{eq:chi}
 \chi(\gamma)=(-1)^{b+c}.
\end{equation}
This is a character. Indeed $a,d$ are odd, and multiplication adds
both $b$ and $c$ modulo $2$. It is trivial on $-I$ and takes the value
$-1$ on each of $T,W_8,C$. Thus \eqref{eq:gensigns} and
Lemma~\ref{lem:generators} imply
\begin{equation}\label{eq:charlaw}
 h|\gamma=\chi(\gamma)h\qquad(\gamma\in\Gamma_0(8)).
\end{equation}

For every $\gamma=\left(\begin{smallmatrix}a&b\\32c&d\end{smallmatrix}\right)
\in\Gamma_0(32)$ we have
\[
 D^{-1}\gamma D=\begin{pmatrix}a&2b\\16c&d\end{pmatrix}
 \in\Gamma_0(8),\qquad
 \chi(D^{-1}\gamma D)=(-1)^{2b+2c}=1.
\]
Since $h=2^{k/2}f|D$, equation \eqref{eq:charlaw} proves
$f|\gamma=f$. Only inclusion of the conjugated group in
$\ker\chi$ is needed; these two groups are not asserted to be equal.

\section{Vanishing at every cusp}

The coefficient bound already obtained implies
\[
 f(x+iy)=O(y^{-A-1})\qquad(0<y\leq1),
\]
uniformly in $x$. It follows that $h|\alpha$ has at most polynomial
growth as $y\to\infty$, uniformly for $x$ in a bounded interval,
for every $\alpha\in\mathrm{GL}_2^+(\mathbb Q)$. To see this when
the lower-left entry of $\alpha$ is nonzero, use
$\operatorname{Im}(\alpha z)=\det(\alpha)y/|cz+d|^2\asymp y^{-1}$
and the factor $(cz+d)^{-k}$. When the lower-left entry is zero,
the Fourier series instead gives exponential decay.

The cusps of $\Gamma_0(8)$ are represented by
$\infty,0,1/2,1/4$, with respective widths $1,8,2,1$.
For a direct completeness check, the four denominators have distinct
greatest common divisors with $8$, so the cusps are inequivalent;
their widths sum to $12=[\SL_2(\mathbb Z):\Gamma_0(8)]$.
They therefore exhaust the orbits of the parabolic action on the
twelve cosets. Representatives and their primitive stabilizers are:
\begin{center}
\renewcommand{\arraystretch}{1.7}
\begin{tabular}{@{}cclc@{}}
\toprule
Cusp & Width $w$ & $P=\alpha T^w\alpha^{-1}$ & $\chi(P)$\\
\midrule
$\infty$ & 1 & $T$ & $-1$\\
$0$ & 8 & $W_8^{-1}$ & $-1$\\
$1/2$ & 2 & $\begin{pmatrix}-3&2\\-8&5\end{pmatrix}$ & $-1$\\[3pt]
$1/4$ & 1 & $\begin{pmatrix}-3&1\\-16&5\end{pmatrix}$ & $-1$\\[3pt]
\bottomrule
\end{tabular}
\end{center}
Here take $\alpha=I,S,
\left(\begin{smallmatrix}1&0\\2&1\end{smallmatrix}\right),
\left(\begin{smallmatrix}1&0\\4&1\end{smallmatrix}\right)$,
respectively. Equation \eqref{eq:charlaw} shows that
\[
 (h|\alpha)(z+w)=-(h|\alpha)(z).
\]
Consequently $h|\alpha$ has a Fourier expansion of period $2w$
containing only odd indices. Polynomial growth excludes every negative
index: equivalently, the Laurent series in $q=e^{2\pi iz/(2w)}$
has no negative powers, by Cauchy's coefficient estimate and the
bound by a power of $\log(1/|q|)$. Anti-periodicity excludes the
constant term. Thus $h|\alpha$ decays exponentially at every one
of these cusps.

Every rational cusp is equivalent to one of them under $\Gamma_0(8)$.
The character law therefore gives vanishing at every rational cusp,
also when $h$ is viewed on the subgroup $D^{-1}\Gamma_0(32)D$.
Rational rescaling preserves this vanishing. Hence $f$ vanishes at
every cusp of $\Gamma_0(32)$, completing the proof of
Theorem~\ref{thm:32}.

\section{What has and has not been generalized}

The translation-module calculation now has the arbitrary parameter
$\tau$ of Proposition~\ref{prop:module}. The defect argument is the
signed criterion of Lemma~\ref{lem:pair}, which is independent of the
level. The half-translation reduction works for every $N$ divisible
by $4$, and its $m=3$ matrix identities work whenever $3\nmid N$.
The integer-conjugacy step closes uniformly for $N=4,8,16,32$.
Among these, level $32$ supplies the additional converse theorem.

This argument does not establish the converse theorem at arbitrary
$4M$. For example, if $M=7$, the parameter in \eqref{eq:r} is $2/7$,
and the known anti-periodicity under $T$ does not control
$\bb(2/7)$. The preceding proofs do not replace it by a scalar.
Further levels require another relation or a different factorization.

The accompanying program \texttt{verify\_N32.py} checks the rational
matrix identities, the parameterized three-dimensional identities,
the signed characters, the four cusp orbits and widths, and the
elliptic trace. It also verifies the recorded SHA-256 hashes of the
three original level $36$ deliverables when they are present next to
this new directory. The mathematical proofs of generation, elliptic
vanishing, and cusp decay are given above and do not depend on a
finite numerical search.

\begin{thebibliography}{9}
\bibitem{CF}
J. B. Conrey and D. W. Farmer,
\emph{An extension of Hecke's converse theorem},
International Mathematics Research Notices 1995, no.~9, 445--463.
\href{https://arxiv.org/abs/math/9502209}{arXiv:math/9502209}.

\bibitem{CFOS}
J. B. Conrey, D. W. Farmer, B. E. Odgers, and N. C. Snaith,
\emph{A converse theorem for $\Gamma_0(13)$},
Journal of Number Theory 122 (2007), 314--323.
\href{https://arxiv.org/abs/math/0601549}{arXiv:math/0601549}.

\bibitem{BCOZ}
B. Bedert, G. Cooper, T. Oliver, and P. Zhang,
\emph{Twisting moduli for $\mathrm{GL}(2)$},
Journal of Number Theory 217 (2020), 142--162,
especially Theorem~2.4 and Section~3.4.
\href{https://arxiv.org/abs/2003.02557}{arXiv:2003.02557}.
\end{thebibliography}
\end{document}
